\chapter{Optical communications}\label{ch:optical}

\file{phy/optical.py} models fiber-optic and free-space optical links, from the laser to the
coherent receiver's digital signal processing. \Cref{fig:optical-chains} shows the two families
of systems it supports.

\begin{figure}[htbp]
\centering
\begin{tikzpicture}[node distance=3mm and 3.6mm,every node/.style={font=\sffamily\scriptsize}]
  \node[lbl,anchor=east] at (-0.15,0) {IM/DD};
  \node[blockbox,minimum width=10mm] (a1) at (0.55,0) {PAM/NRZ};
  \node[blockbox,right=of a1,minimum width=10mm] (a2) {laser /\\EML};
  \node[warmbox,right=of a2,minimum width=10mm] (a3) {fiber};
  \node[warmbox,right=of a3,minimum width=10mm] (a4) {EDFA};
  \node[greenbox,right=of a4,minimum width=10mm] (a5) {PIN /\\APD};
  \node[greenbox,right=of a5,minimum width=10mm] (a6) {AGC};
  \node[greenbox,right=of a6,minimum width=10mm] (a7) {decision};
  \foreach \x/\y in {a1/a2,a2/a3,a3/a4,a4/a5,a5/a6,a6/a7} \draw[sig] (\x)--(\y);
  \node[lbl,anchor=east] at (-0.15,-1.35) {coherent};
  \node[blockbox,minimum width=10mm] (b1) at (0.55,-1.35) {DP-QAM};
  \node[blockbox,right=of b1,minimum width=10mm] (b2) {IQ-MZM};
  \node[warmbox,right=of b2,minimum width=12mm] (b3) {$N\times$(span\\+EDFA)};
  \node[greenbox,right=of b3,minimum width=10mm] (b4) {coherent\\front end};
  \node[greenbox,right=of b4,minimum width=8mm] (b5) {CD};
  \node[greenbox,right=of b5,minimum width=8mm] (b6) {FO};
  \node[greenbox,right=of b6,minimum width=10mm] (b7) {$2{\times}2$\\MIMO};
  \node[greenbox,right=of b7,minimum width=8mm] (b8) {CPR};
  \foreach \x/\y in {b1/b2,b2/b3,b3/b4,b4/b5,b5/b6,b6/b7,b7/b8} \draw[sig] (\x)--(\y);
\end{tikzpicture}
\caption{Intensity modulation with direct detection, and dual-polarization coherent transmission.}
\label{fig:optical-chains}
\end{figure}

\begin{incode}
Optical signals are complex field envelopes around the optical carrier, scaled so that
$|E|^2$ is the instantaneous optical power in watts. Dual polarization uses the shape
$(B,2,n)$. Optical blocks maintain the following metadata:
\begin{itemize}
  \item \meta{wavelength_nm};
  \item \meta{power}, the average launch or received power in W;
  \item \meta{cd_ps_nm}, the accumulated chromatic dispersion $D\ell$;
  \item \meta{ase_psd}, the accumulated ASE density per polarization in W/Hz.
\end{itemize}
The OSNR is computed from the last two at any point of the link.
\end{incode}

% ==================================================================
\section{Transmitters}
\subsection{Directly modulated lasers}
An intensity modulator maps the electrical drive $x\in[-1,1]$ to the optical power
\begin{equation}
  P(t)=P_{\mathrm{avg}}\big(1+m\,x(t)\big),\qquad m=\frac{\mathrm{ER}-1}{\mathrm{ER}+1},
\end{equation}
since the extinction ratio is $\mathrm{ER}=P_1/P_0=(1+m)/(1-m)$. Three non-idealities are
modelled:
\begin{itemize}
  \item \textbf{Relative intensity noise}: multiplicative Gaussian noise of variance
  $\mathrm{RIN}\cdot f_s/2$ on $P$.
  \item \textbf{Phase noise}: a Wiener process with the laser linewidth (\cref{eq:wiener}).
  \item \textbf{Transient chirp} of a directly modulated laser. The instantaneous frequency
  follows the power changes, $\Delta\nu(t)=\frac{\alpha}{4\pi}\frac{d}{dt}\ln P(t)$, with $\alpha$
  the Henry linewidth-enhancement factor. Integrating gives the phase
  \begin{equation}
    \phi_{\mathrm{chirp}}(t)=\frac{\alpha}{2}\ln\frac{P(t)}{P_{\mathrm{avg}}} .
  \end{equation}
\end{itemize}
The field is $E=\sqrt P\,e^{j\phi}$.
\codeexcerpt{im_mod}{Laser intensity modulation with RIN, phase noise and chirp (\blk{intensity_modulator}).}

\begin{pitfall}[Not a rate-equation model]
The laser is modelled behaviourally (power, chirp, linewidth). Relaxation oscillations, the
adiabatic chirp term and thermal effects, which require the laser rate equations, are not
included.
\end{pitfall}

\subsection{Mach--Zehnder and IQ modulators}
A Mach--Zehnder modulator in push--pull configuration, biased at null, has the field transfer
function $E_{\mathrm{out}}/E_{\mathrm{in}}=\sin\big(\frac{\pi}{2}V/V_\pi\big)$. An IQ modulator
nests two of them, with a $90^\circ$ phase shift between the arms:
\begin{equation}
  \frac{E_{\mathrm{out}}}{E_{\mathrm{in}}}\propto\sin\Big(\frac\pi2\frac{V_I}{V_\pi}\Big)+j\sin\Big(\frac\pi2\frac{V_Q}{V_\pi}\Big).
\end{equation}
With the peak drive set to \prm{drive_swing}$\times V_\pi$, small swings are nearly linear,
while large swings compress the outer constellation points (\cref{fig:mzm}). The \prm{linear}
model maps the drive directly onto the field. The laser power is split equally between the
polarizations.

\begin{figure}[htbp]
\centering
\begin{tikzpicture}
\begin{axis}[width=9cm,height=5cm,xlabel={drive $V/V_\pi$},ylabel={normalized output},xmin=-1,xmax=1,
  ymin=-1.1,ymax=1.1,legend pos=north west,samples=101]
  \addplot[th=accent,domain=-1:1] {sin(90*x)};
  \addplot[th=accent2,domain=-1:1] {sin(90*x)^2};
  \addplot[th=black!50,dashed,domain=-0.6:0.6] {pi/2*x};
  \draw[accent3,fill=accent3!15,opacity=0.5] (axis cs:-0.5,-1.1) rectangle (axis cs:0.5,1.1);
  \legend{field $\sin(\frac\pi2 V/V_\pi)$,power,linear approx.}
\end{axis}
\end{tikzpicture}
\caption{MZM transfer function at null bias. The shaded region is the default drive swing
($0.5\,V_\pi$).}
\label{fig:mzm}
\end{figure}
\codeexcerpt{iq_mzm}{Sinusoidal MZM transfer and power normalization (\blk{iq_modulator}).}

% ==================================================================
\section{Propagation in single-mode fiber}\label{sec:nlse}
\subsection{The nonlinear Schrödinger equation}
The slowly varying envelope $A(z,t)$ of a single-mode fiber, in a frame moving with the group
velocity, obeys
\begin{equation}
  \frac{\partial A}{\partial z}=
  \underbrace{-\frac{\alpha}{2}A-j\frac{\beta_2}{2}\frac{\partial^2A}{\partial t^2}}_{\hat D A}
  \;+\;\underbrace{j\gamma|A|^2A}_{\hat N A},
  \label{eq:nlse}
\end{equation}
with three ingredients~\cite{agrawal2019}:
\begin{itemize}
  \item the power attenuation $\alpha$, given in dB/km as \prm{attenuation_db_km};
  \item the group-velocity dispersion $\beta_2$, related to the dispersion parameter $D$
  (ps/nm/km) by
  \begin{equation}
    \beta_2=-\frac{D\lambda^2}{2\pi c}\qquad(\beta_2\approx\SI{-21.7}{ps^2/km}\ \text{for }D=17,\ \lambda=1550\,\mathrm{nm});
  \end{equation}
  \item the Kerr coefficient $\gamma$, given in 1/W/km as \prm{gamma_w_km}.
\end{itemize}
With the Fourier convention of \cref{eq:fourier-convention}, $\partial_t\to j\omega$, so the
linear part has the exact solution
\begin{equation}
  \tilde A(z,\omega)=\tilde A(0,\omega)\exp\!\Big(j\frac{\beta_2}{2}\omega^2z-\frac{\alpha}{2}z\Big),
  \label{eq:dispersion-operator}
\end{equation}
an all-pass filter with quadratic phase. The dispersion therefore depends only on the
accumulated $D\ell$ (ps/nm), the quantity tracked in \meta{cd_ps_nm}.
\codeexcerpt{beta2}{From $D$ to $\beta_2$.}
\codeexcerpt{disp_op}{The dispersion operator of \cref{eq:dispersion-operator}.}

\paragraph{Dispersion alone.}
An unchirped Gaussian pulse $A(0,t)=e^{-t^2/(2T_0^2)}$ stays Gaussian and broadens as
\begin{equation}
  \frac{T_1(z)}{T_0}=\sqrt{1+\Big(\frac{z}{L_D}\Big)^2},\qquad L_D=\frac{T_0^2}{|\beta_2|}.
\end{equation}
For $T_0=5$\,ps and $|\beta_2|=\fiberbeta$\,ps$^2$/km, $L_D=\fiberLD$\,km.
\Cref{fig:dispersion} compares the \code{propagate} solver with this result. The simulated RMS
widths agree with the formula to all printed digits, as expected from an exact spectral method.

\begin{figure}[htbp]
\centering
\begin{tikzpicture}
\begin{groupplot}[group style={group size=2 by 1,horizontal sep=16mm},width=6.9cm,height=5.2cm]
\nextgroupplot[xlabel={$t$ [ps]},ylabel={$|A|^2/P_0$},xmin=-40,xmax=40,ymin=0,ymax=1.05,
  legend pos=north east,title={pulse shape}]
  \addplot[th=black!70] table[x=t_ps,y=z0]{\datafile{fiber_pulses}};
  \addplot[th=accent] table[x=t_ps,y=z1]{\datafile{fiber_pulses}};
  \addplot[th=accent3] table[x=t_ps,y=z2]{\datafile{fiber_pulses}};
  \addplot[th=accent2] table[x=t_ps,y=z4]{\datafile{fiber_pulses}};
  \legend{$z=0$,$L_D$,$2L_D$,$4L_D$}
\nextgroupplot[xlabel={$z/L_D$},ylabel={$T_1/T_0$},xmin=0,xmax=5,ymin=0.8,ymax=5.4,legend pos=north west,
  title={RMS broadening}]
  \addplot[th=accent] table[x=z,y=theory]{\datafile{fiber_broadening}};
  \addplot[sim=accent2] table[x=z,y=sim]{\datafile{fiber_broadening}};
  \legend{$\sqrt{1+(z/L_D)^2}$,\code{propagate}}
\end{groupplot}
\end{tikzpicture}
\caption{Chromatic dispersion of a 5-ps Gaussian pulse in standard single-mode fiber
($D=17$\,ps/nm/km, no loss).}
\label{fig:dispersion}
\end{figure}

\paragraph{Kerr nonlinearity alone.}
Without dispersion \cref{eq:nlse} integrates exactly. The power profile is unchanged and the phase
grows in proportion to it, $A(L,t)=A(0,t)\,e^{-\alpha L/2}e^{j\gamma|A(0,t)|^2L_{\mathrm{eff}}}$ with
$L_{\mathrm{eff}}=(1-e^{-\alpha L})/\alpha$. This self-phase modulation (SPM) chirps the pulse and
broadens its spectrum, with a number of spectral peaks that grows with the peak phase
$\phi_{\max}=\gamma P_0L_{\mathrm{eff}}$ (\cref{fig:spm}).

\begin{figure}[htbp]
\centering
\begin{tikzpicture}
\begin{groupplot}[group style={group size=3 by 1,horizontal sep=4mm,y descriptions at=edge left},
  width=5.1cm,height=4.4cm,xmin=-1.5,xmax=1.5,ymin=0,ymax=1.05,xlabel={$(\nu-\nu_0)T_0$}]
\nextgroupplot[ylabel={spectral density},title={$\phi_{\max}=0$}]
  \addplot[th=black!70] table[x=f,y=p0]{\datafile{spm_spectra}};
\nextgroupplot[title={$\phi_{\max}=1.5\pi$}]
  \addplot[th=accent] table[x=f,y=p15]{\datafile{spm_spectra}};
\nextgroupplot[title={$\phi_{\max}=3.5\pi$}]
  \addplot[th=accent2] table[x=f,y=p35]{\datafile{spm_spectra}};
\end{groupplot}
\end{tikzpicture}
\caption{SPM-broadened spectra of a Gaussian pulse computed by \code{propagate} with $D=0$,
the classic figure of Agrawal~\cite{agrawal2019}: $\phi_{\max}=(M-\frac12)\pi$ produces $M$ peaks.}
\label{fig:spm}
\end{figure}

\subsection{The split-step Fourier method}
With both terms present, \cref{eq:nlse} has no general closed form. Formally
$A(z+h)=e^{h(\hat D+\hat N)}A(z)$. The two operators do not commute, but each is trivial to
apply on its own: $\hat D$ is a multiplication in the frequency domain, and $\hat N$ a
multiplication in the time domain. The symmetric (Strang) splitting
\begin{equation}
  A(z+h)=e^{\frac h2\hat D}\,e^{h\hat N}\,e^{\frac h2\hat D}A(z)+\mathcal O(h^3)
  \label{eq:strang}
\end{equation}
is exact to second order, by the Baker--Campbell--Hausdorff formula: the first commutator
$[\hat D,\hat N]$ cancels by symmetry. Over a fixed length the $L/h$ local errors add up to a
global error $\mathcal O(h^2)$.

Attenuation is included in the half steps. The nonlinear phase is applied at mid-step to the
field whose power has been reduced by the first half step. The effective length that makes
this exact for a constant envelope is
\begin{equation}
  L_{\mathrm{eff}}(h)=\int_{-h/2}^{h/2}e^{-\alpha z'}dz'=\frac{2\sinh(\alpha h/2)}{\alpha} .
\end{equation}
For dual polarization the Manakov equation replaces $|A|^2$ by $|A_x|^2+|A_y|^2$ and $\gamma$ by
$\frac89\gamma$, the average over fast random polarization evolution.
\codeexcerpt{ssfm}{The split-step Fourier solver of the NLSE / Manakov equation (\code{propagate}).}

\begin{figure}[htbp]
\centering
\begin{tikzpicture}
\begin{groupplot}[group style={group size=2 by 1,horizontal sep=17mm},width=6.9cm,height=5.4cm]
\nextgroupplot[xmode=log,ymode=log,xlabel={step $h$ [km]},ylabel={relative error after $4L_D$},
  xmin=0.005,xmax=4,ymin=1e-6,ymax=1,legend pos=north west,title={convergence}]
  \addplot[sim=accent2,mark size=2pt] table[x=h,y=err]{\datafile{soliton_convergence}};
  \addplot[th=black!50,dashed,domain=0.005:3] {0.09*x^2};
  \legend{split-step,slope 2 ($\propto h^2$)}
  \draw[accent,densely dotted] (axis cs:2,1e-6) -- (axis cs:2,1) node[pos=0.5,left,font=\scriptsize,rotate=90,anchor=south]{default 0.5 step/km};
\nextgroupplot[xlabel={$t$ [ps]},ylabel={$|A|^2/P_0$},xmin=-30,xmax=30,ymin=0,ymax=1.05,
  legend pos=north east,title={after $4L_D$}]
  \addplot[th=black!70] table[x=t_ps,y=input]{\datafile{soliton_profiles}};
  \addplot[th=accent,dashed] table[x=t_ps,y=linear]{\datafile{soliton_profiles}};
  \addplot[sim={accent2,mark size=0.9pt}] table[x=t_ps,y=nlse]{\datafile{soliton_profiles}};
  \legend{input / exact,dispersion only,NLSE (split-step)}
\end{groupplot}
\end{tikzpicture}
\caption{Validation of the solver on the fundamental soliton, an exact solution of
\cref{eq:nlse} with $\alpha=0$. With $T_0=5$\,ps, $P_0=|\beta_2|/(\gamma T_0^2)=\solitonPzero$\,W
and $L=\solitonL$\,km $=4L_D$, the pulse must return unchanged. Left: the relative $L^2$ error
falls as $h^2$, confirming the order of \cref{eq:strang}. Right: dispersion alone would broaden
the pulse, while the balance with SPM preserves it.}
\label{fig:soliton}
\end{figure}

\begin{pitfall}[Choosing the step]
The default \prm{steps_per_km}\code{=0.5} (2-km steps) is adequate for typical launch powers
around 0\,dBm per channel, where the nonlinear length $1/(\gamma P)$ is hundreds of km. The
soliton of \cref{fig:soliton} carries 28\,dBm, a nonlinear length of about 1\,km. There the
default step leaves a 12\,\% error, and 16 steps per km are needed for $10^{-3}$. For high powers,
dense WDM or solitons, reduce the step until the result stops changing. Without \prm{nonlinear} the solver applies one exact
FFT for the whole span.
\end{pitfall}

\subsection{Polarization-mode dispersion}
Random birefringence rotates the state of polarization and delays the two principal states by
the differential group delay (DGD). The first-order model is
\begin{equation}
  \tilde{\vect A}_{\mathrm{out}}(\omega)=\mat U_2\begin{pmatrix}e^{j\omega\tau/2}&0\\0&e^{-j\omega\tau/2}\end{pmatrix}\mat U_1\tilde{\vect A}_{\mathrm{in}}(\omega),
  \qquad\tau=D_{\mathrm{PMD}}\sqrt{\ell},
\end{equation}
with $\mat U_{1,2}$ Haar-random unitary matrices. These are generated by a QR decomposition of
a complex Gaussian matrix with the phases of the diagonal of $\mat R$ removed.
\codeexcerpt{pmd}{Random polarization rotation and first-order PMD.}
Here $\tau$ is set to its mean value; a full statistical model would draw it from a Maxwellian
distribution.

% ==================================================================
\section{Optical amplification and OSNR}
\subsection{ASE noise}
An erbium-doped fiber amplifier of gain $G$ adds amplified spontaneous emission with power
spectral density, per polarization,
\begin{equation}
  S_{\mathrm{ASE}}=n_{\mathrm{sp}}\,h\nu\,(G-1),\qquad \mathrm{NF}=\frac{2n_{\mathrm{sp}}(G-1)}{G}\ \Rightarrow\ n_{\mathrm{sp}}=\frac{\mathrm{NF}\cdot G}{2(G-1)},
\end{equation}
where the noise figure NF is at least 2 (3\,dB) for $n_{\mathrm{sp}}\ge1$. In the simulation band
$f_s$, the ASE is circular Gaussian noise of variance $S_{\mathrm{ASE}}f_s$ per polarization.
\codeexcerpt{ase}{ASE density from gain and noise figure.}
\codeexcerpt{add_ase}{ASE noise in the simulated bandwidth.}

\subsection{Optical signal-to-noise ratio}
The OSNR is measured in a reference bandwidth $B_{\mathrm{ref}}=12.5$\,GHz (0.1\,nm at 1550\,nm),
counting the noise in both polarizations:
\begin{equation}
  \mathrm{OSNR}=\frac{P}{2S_{\mathrm{ASE}}B_{\mathrm{ref}}} .
\end{equation}
In a chain of $N$ identical spans, each followed by an amplifier compensating its loss
($G=e^{\alpha L_{\mathrm{span}}}$), the ASE adds up while the signal power is restored. With the
NF definition above, $2S_{\mathrm{ASE}}\approx\mathrm{NF}\,G\,h\nu$, so
\begin{equation}
  \mathrm{OSNR}_{\mathrm{dB}}\approx\underbrace{-10\log_{10}\frac{h\nu B_{\mathrm{ref}}}{1\,\mathrm{mW}}}_{58.0\ \mathrm{dB}}
  +P_{\mathrm{dBm}}-G_{\mathrm{dB}}-\mathrm{NF}_{\mathrm{dB}}-10\log_{10}N,
  \label{eq:osnr58}
\end{equation}
the ``58\,dB rule'' of link design.
\codeexcerpt{link_loop}{The multi-span link: propagate, amplify, accumulate the ASE density.}

\begin{figure}[htbp]
\centering
\begin{tikzpicture}
\begin{semilogxaxis}[width=10cm,height=5.4cm,xlabel={number of spans $N$},ylabel={OSNR [dB]},
  xmin=0.8,xmax=50,ymin=18,ymax=40,legend pos=north east]
  \addplot[th=accent] table[x=spans,y=rule58]{\datafile{osnr_spans}};
  \addplot[sim=accent2] table[x=spans,y=sim]{\datafile{osnr_spans}};
  \legend{58\,dB rule \eqref{eq:osnr58},\blk{optical_link} metric}
\end{semilogxaxis}
\end{tikzpicture}
\caption{OSNR after $N\times80$\,km (0.2\,dB/km, NF\,=\,5\,dB, 0\,dBm per channel). The block's
metric and the rule differ by less than 0.05\,dB, the effect of the $G-1\approx G$ approximation
and of the exact photon energy.}
\label{fig:osnr}
\end{figure}

\subsection{From OSNR to $E_s/N_0$}
For a coherent receiver whose noise is the ASE, the noise density per polarization is
$N_0=S_{\mathrm{ASE}}$ and the signal power per polarization is $P/p$ for $p$ polarizations.
Hence
\begin{equation}
  \frac{E_s}{N_0}=\frac{P/p}{S_{\mathrm{ASE}}R_s}=\mathrm{OSNR}\cdot\frac{2B_{\mathrm{ref}}}{p\,R_s},
\end{equation}
implemented as \code{osnr_to_esn0_db}. The test-suite uses it to check that the full coherent
DP-QPSK link (\cref{sec:coherent-dsp}) performs within about 0.3\,dB of theory.

% ==================================================================
\section{Wavelength-division multiplexing}
\blk{wdm_multiplexer} surrounds the channel under test with $N_{\mathrm{ch}}-1$ neighbours at
multiples of the channel spacing. The neighbours carry independent symbols with the same
constellation, pulse shape and power. Propagating the composite field through a nonlinear
fiber produces cross-phase modulation (the power of the neighbours modulates the phase of the
channel under test) and four-wave mixing. \blk{wdm_demultiplexer} selects a channel with a
second-order super-Gaussian optical filter, $|H(f)|=\exp\!\big(-\tfrac{\ln2}{2}(2f/B)^4\big)$, which
is 3\,dB down at $\pm B/2$. The sample rate must cover all channels. The \emph{WDM with fiber
nonlinearity} scenario shows the MER of the centre channel dropping by several dB when the
neighbours are switched on.
\codeexcerpt{wdm_demux}{Channel selection by frequency shift and optical filtering.}

% ==================================================================
\section{Free-space optics}
Atmospheric turbulence makes the received intensity fluctuate (scintillation). \blk{fso_turbulence}
multiplies the field by $\sqrt h$, where the normalized irradiance $h$ ($\E h=1$) is constant
over a frame and follows one of two models.
\begin{itemize}
  \item \textbf{Log-normal} (weak turbulence): $\ln h\sim\N(-\sigma^2/2,\sigma^2)$.
  \item \textbf{Gamma-gamma} (moderate to strong turbulence): $h$ is the product of two
  independent unit-mean gamma variables with parameters $\alpha$ (large eddies) and $\beta$
  (small eddies). Its density is
  \begin{equation}
    f(h)=\frac{2(\alpha\beta)^{\frac{\alpha+\beta}{2}}}{\Gamma(\alpha)\Gamma(\beta)}\,h^{\frac{\alpha+\beta}{2}-1}K_{\alpha-\beta}\big(2\sqrt{\alpha\beta h}\big).
  \end{equation}
\end{itemize}
\codeexcerpt{fso}{Scintillation models.}

\begin{figure}[htbp]
\centering
\begin{tikzpicture}
\begin{axis}[width=10cm,height=5.2cm,xlabel={normalized irradiance $h$},ylabel={pdf},xmin=0,xmax=3,ymin=0,
  legend pos=north east]
  \addplot[th=accent] table[x=h,y=ln_pdf]{\datafile{fso_pdf}};
  \addplot[sim=accent,mark size=1.2pt] table[x=h,y=ln_hist]{\datafile{fso_pdf}};
  \addplot[th=accent2] table[x=h,y=gg_pdf]{\datafile{fso_pdf}};
  \addplot[sim={accent2,mark=square*},mark size=1.1pt] table[x=h,y=gg_hist]{\datafile{fso_pdf}};
  \legend{log-normal $\sigma^2{=}0.1$,histogram,{gamma-gamma $\alpha{=}4$, $\beta{=}2$},histogram}
\end{axis}
\end{tikzpicture}
\caption{Histograms of $h$ from \blk{fso_turbulence} (60\,000 frames) against the closed-form densities.}
\label{fig:fso}
\end{figure}

% ==================================================================
\section{Direct detection}
\subsection{PIN and APD receivers}
A photodiode converts power into current, $I=\mathcal RP$, with responsivity $\mathcal R$
(A/W). Over the noise bandwidth $B=f_s/2$ the current carries three contributions:
\begin{itemize}
  \item \textbf{shot noise}, $\sigma_{\mathrm{sh}}^2=2qM^2F\,(\mathcal RP+I_d)B$;
  \item \textbf{thermal noise}, $\sigma_{\mathrm{th}}^2=i_n^2B$, with $i_n$ the input-referred
  current density (\prm{thermal_noise_pa_rthz});
  \item \textbf{avalanche gain} (APD): the multiplication $M$ comes with McIntyre's excess noise
  factor $F=kM+(1-k)(2-1/M)$, with $k$ the ionization ratio.
\end{itemize}
The APD gain raises the signal above thermal noise but amplifies shot noise by $M^2F$, so there
is an optimum $M$.
\codeexcerpt{photodiode}{Square-law detection with shot, thermal and APD excess noise.}

\subsection{OOK sensitivity and the Q factor}
For on--off keying with levels $I_{1,0}=\mathcal RP_{1,0}$ and Gaussian noise of standard
deviations $\sigma_{1,0}$, the optimum threshold gives
\begin{equation}
  \mathrm{BER}=\frac12\erfc\Big(\frac{Q}{\sqrt2}\Big),\qquad Q=\frac{I_1-I_0}{\sigma_1+\sigma_0},
  \qquad P_{1,0}=\frac{2P_{\mathrm{avg}}}{1+\mathrm{ER}^{\mp1}} .
\end{equation}
Thermal noise is signal independent, so $\sigma_1=\sigma_0$ and $Q\propto P_{\mathrm{avg}}$: the BER
falls by many decades within a few dB (\cref{fig:ook}). The receiver sensitivity is
conventionally quoted at $10^{-9}$ ($Q=6$) or at a FEC threshold.

\begin{figure}[htbp]
\centering
\begin{tikzpicture}
\begin{semilogyaxis}[width=10cm,height=6cm,xlabel={average received power [dBm]},ylabel={BER},
  xmin=-28,xmax=-18,ymin=1e-7,ymax=0.5,legend pos=south west]
  \addplot[th=accent] table[x=p,y=ber]{\datafile{ook_theory}};
  \addplot[sim=accent2] table[x=p,y=ber]{\datafile{ook_sim}};
  \legend{Q-factor theory,simulated}
\end{semilogyaxis}
\end{tikzpicture}
\caption{10\,Gb/s NRZ-OOK into a thermal-noise-limited PIN receiver ($\mathcal R=0.8$\,A/W,
20\,pA/$\sqrt{\mathrm{Hz}}$, ER\,=\,10\,dB): chain \blk{intensity_modulator}$\to$\blk{photodiode}$\to$\blk{agc}$\to$\blk{line_decoder}
against the Q-factor formula.}
\label{fig:ook}
\end{figure}

\begin{incode}
\blk{agc} removes the mean (the DC photocurrent) and scales each frame to the nominal symbol
power of the upstream line code or constellation. It subtracts the known noise variance when
estimating the signal power, so the decision thresholds stay correct at low SNR.
\end{incode}

% ==================================================================
\section{Coherent detection}
\subsection{The coherent front end}
A coherent receiver mixes the signal with a local oscillator (LO) in a $90^\circ$ hybrid and
detects the four outputs with two balanced photodiode pairs per polarization
(\cref{fig:coherent-fe}). Balanced detection cancels the direct-detection terms $|E|^2$ and
$|E_{\mathrm{LO}}|^2$ and leaves the beat term:
\begin{equation}
  i_I+j\,i_Q=\mathcal R\sqrt{P_{\mathrm{LO}}}\;E(t)\,e^{-j(2\pi\Delta f\,t+\phi_{\mathrm{LO}}(t))}+n(t).
\end{equation}
The field is thus translated linearly to the electrical domain, including its phase and
polarization. With a strong LO, the LO shot noise $2q\mathcal RP_{\mathrm{LO}}B$ dominates the
thermal noise.

\begin{figure}[htbp]
\centering
\begin{tikzpicture}[node distance=5mm and 8mm]
  \node[blockbox] (pbs) {PBS};
  \node[blockbox,right=of pbs,yshift=10mm,minimum height=11mm] (hx) {$90^\circ$\\hybrid X};
  \node[blockbox,right=of pbs,yshift=-10mm,minimum height=11mm] (hy) {$90^\circ$\\hybrid Y};
  \node[warmbox,above=4mm of hx] (lox) {LO ($\Delta f$, $\Delta\nu$)};
  \node[warmbox,below=4mm of hy] (loy) {LO ($\Delta f$, $\Delta\nu$)};
  \node[greenbox,right=of hx,minimum width=13mm] (bx) {2 balanced\\PDs};
  \node[greenbox,right=of hy,minimum width=13mm] (by) {2 balanced\\PDs};
  \node[blockbox,right=of bx] (ax) {ADC};
  \node[blockbox,right=of by] (ay) {ADC};
  \node[left=6mm of pbs,lbl] (in) {$E_x,E_y$};
  \draw[sig] (in)--(pbs); \draw[sig] (pbs.east)--(hx.west); \draw[sig] (pbs.east)--(hy.west);
  \draw[sig] (lox)--(hx); \draw[sig] (loy)--(hy);
  \draw[sig] (hx)--(bx); \draw[sig] (hy)--(by); \draw[sig] (bx)--node[above,lbl]{$I_x,Q_x$}(ax); \draw[sig] (by)--node[above,lbl]{$I_y,Q_y$}(ay);
\end{tikzpicture}
\caption{Polarization-diverse coherent front end, modelled by \blk{coherent_receiver}. One LO laser
is split between the two hybrids.}
\label{fig:coherent-fe}
\end{figure}

\begin{incode}
The simulator processes each frame with circular FFTs (fiber, dispersion compensation). An LO
phase that does not return to its initial value at the end of the frame would create a
discontinuity at the frame boundary, spread over the frame by the dispersion compensator. The
block therefore makes the LO phase periodic:
\begin{itemize}
  \item the frequency offset is rounded to the frame's frequency grid $f_s/n$;
  \item the Wiener phase is turned into a Brownian bridge, $\phi_k-\frac kn\phi_n$, which has the
  same local increments but returns to zero.
\end{itemize}
\end{incode}
\codeexcerpt{coherent_rx}{Coherent detection with a periodic LO phase (\blk{coherent_receiver}).}

\subsection{Coherent digital signal processing}\label{sec:coherent-dsp}
The chain follows the architecture described by Savory~\cite{savory2010}.
\Cref{fig:coherent-const} follows a 128\,Gb/s DP-QPSK signal through the DSP chain after
$20\times80$\,km, with random polarization, a 1\,GHz LO offset and 100\,kHz lasers.

\begin{figure}[htbp]
\centering
\begin{tikzpicture}
\begin{groupplot}[group style={group size=4 by 1,horizontal sep=4mm,y descriptions at=edge left},
  constaxis,width=4.1cm,height=4.1cm,xmin=-2.2,xmax=2.2,ymin=-2.2,ymax=2.2,xtick={-2,0,2},ytick={-2,0,2}]
\nextgroupplot[title={front-end output}]
  \addplot[scatterpts=black!50] table[x=i,y=q]{\datafile{coh_rx}};
\nextgroupplot[title={after CD + FO comp.}]
  \addplot[scatterpts=accent4] table[x=i,y=q]{\datafile{coh_fo}};
\nextgroupplot[title={after $2{\times}2$ equalizer}]
  \addplot[scatterpts=accent] table[x=i,y=q]{\datafile{coh_eq}};
\nextgroupplot[title={after phase recovery}]
  \addplot[scatterpts=accent3] table[x=i,y=q]{\datafile{coh_cpr}};
\end{groupplot}
\end{tikzpicture}
\caption{Polarization X of the \emph{Coherent DP-QPSK long haul} scenario at four points of
the DSP chain (normalized to unit power).}
\label{fig:coherent-const}
\end{figure}

\paragraph{Chromatic dispersion compensation.}
The accumulated dispersion is known from the metadata (or given), and its all-pass response
\cref{eq:dispersion-operator} is inverted exactly by multiplying with its conjugate in the
frequency domain.
\codeexcerpt{cd_comp}{Frequency-domain CD compensation.}

\paragraph{Frequency-offset estimation.}
Raising QPSK to the fourth power removes the modulation: $(s\,e^{j2\pi\Delta f t})^4=s^4e^{j2\pi4\Delta ft}$
with $s^4=-1$ for every QPSK symbol, and approximately constant on average for square QAM. The
spectrum of $r^4$ therefore has a line at $4\Delta f$. The estimate is the position of its peak
divided by 4, refined by zero-padding and computed on a low-pass filtered copy to reject noise.
In dual polarization a rotated mixture $as_1+bs_2$ has a fourth-power line proportional to
$a^4+b^4$, which can vanish. The estimator therefore also uses the combinations
$(x_0+x_1)/\sqrt2$ and $(x_0+jx_1)/\sqrt2$.

Compensating the offset \emph{after} the dispersion compensator leaves a residual: the CD
filter was applied to a spectrum shifted by $\Delta\omega$, so the error is
$\exp\!\big(j\frac{\beta_2L}{2}(2\omega\Delta\omega+\Delta\omega^2)\big)$, a bulk delay
$\beta_2L\Delta\omega$ plus a constant phase. The block removes it.
\codeexcerpt{fo_est}{Fourth-power FO estimation over several polarization combinations.}

\paragraph{Adaptive $2\times2$ equalization.}
Polarization rotation, PMD and residual dispersion form a $2\times2$ MIMO channel, inverted by
a fractionally spaced butterfly of four FIR filters (\cref{fig:butterfly}):
\begin{equation}
  y_p[k]=\sum_{q\in\{x,y\}}\vect w_{pq}\T\,\vect u_q[k],\qquad
  \vect w_{pq}\leftarrow\vect w_{pq}+\mu\,e_p[k]\,\vect u_q\conj[k] .
\end{equation}
The error $e_p$ is either:
\begin{itemize}
  \item the constant-modulus error $y_p(R-|y_p|^2)$ with $R=\E|s|^4/\E|s|^2$ (Godard's
  CMA~\cite{godard1980}), which is blind and insensitive to the carrier phase; or
  \item the LMS error $d_p-y_p$, trained on known symbols and then decision-directed, with a
  larger step during training.
\end{itemize}

\begin{figure}[htbp]
\centering
\begin{tikzpicture}[x=1cm,y=1cm]
  % input x feeds w_xx (to x out) and w_yx (to y out); input y feeds w_xy and w_yy
  \node[lbl,anchor=east] (xi) at (0,0) {$x$ in};
  \node[lbl,anchor=east] (yi) at (0,-3.3) {$y$ in};
  \node[blockbox,minimum width=11mm] (wxx) at (2.2,0.2) {$\vect w_{xx}$};
  \node[blockbox,minimum width=11mm] (wyx) at (2.2,-1.0) {$\vect w_{yx}$};
  \node[blockbox,minimum width=11mm] (wxy) at (2.2,-2.3) {$\vect w_{xy}$};
  \node[blockbox,minimum width=11mm] (wyy) at (2.2,-3.5) {$\vect w_{yy}$};
  \node[sumnode] (sx) at (4.4,0) {$+$};
  \node[sumnode] (sy) at (4.4,-3.3) {$+$};
  \node[lbl,anchor=west] (xo) at (5.2,0) {$x$ out};
  \node[lbl,anchor=west] (yo) at (5.2,-3.3) {$y$ out};
  \coordinate (bx) at (0.7,0); \coordinate (by) at (0.7,-3.3);
  \draw[thick] (xi.east) -- (bx); \draw[thick] (yi.east) -- (by);
  \draw[sig] (bx) |- (wxx.west); \draw[sig] (bx) |- (wyx.west);
  \draw[sig] (by) |- (wxy.west); \draw[sig] (by) |- (wyy.west);
  \fill (bx) circle (1.2pt); \fill (by) circle (1.2pt);
  \draw[sig] (wxx.east) -- (sx); \draw[sig] (wxy.east) -- (sx.south west);
  \draw[sig] (wyy.east) -- (sy); \draw[sig] (wyx.east) -- (sy.north west);
  \draw[sig] (sx) -- (xo); \draw[sig] (sy) -- (yo);
\end{tikzpicture}
\caption{The butterfly equalizer (2 samples per symbol, \prm{n_taps} taps per filter).}
\label{fig:butterfly}
\end{figure}
\codeexcerpt{butterfly}{Butterfly update with CMA or trained/decision-directed LMS.}

\paragraph{Carrier phase recovery.}
Laser phase noise of $2\times100$\,kHz at 32\,GBd is slow, but large enough to need
feed-forward tracking. Two estimators are available.
\begin{itemize}
  \item The Viterbi--Viterbi estimator removes the QPSK modulation with the fourth power:
  $\hat\phi=\frac14\big(\arg\sum_{|i|\le W}y_{k+i}^4-\pi\big)$.
  \item Blind phase search (BPS)~\cite{pfau2009} works for any square QAM. It rotates the
  symbols by $B$ test phases in $[-\pi/4,\pi/4)$, sums over a window the squared distances to
  the nearest constellation points, and keeps the phase with the smallest sum.
\end{itemize}
Both estimates are defined modulo $\pi/2$; they are unwrapped (on $4\hat\phi$) to follow the phase
continuously. The remaining quadrant ambiguity is resolved once per frame with a short block of
known symbols, taken after the equalizer has converged.
\codeexcerpt{bps}{Viterbi--Viterbi and blind phase search.}
